\section*{الفصل \ref{ch:g10:freshwater-resources} --- موارد المياه العذبة}

\begin{solution}{exo:g10:freshwater-resources:1}
\omterm{def:g10:freshwater-resources:stock}{المخزون} هو كمية الماء المحفوظة في خزان في لحظة معينة، بالوحدة
\unit{km^3} (أو بالكيلوغرامات). \omterm{def:g10:freshwater-resources:stock}{والدفق} هو الكمية التي تنتقل من خزان إلى آخر في وحدة
الزمن، بالوحدة \unit{km^3} في السنة. \omterm{def:g10:freshwater-resources:residence-time}{وزمن المكوث} هو متوسط الزمن الذي يقضيه الماء في
خزان، بالأيام أو السنوات؛ وهو في حالة الثبات \omterm{def:g10:freshwater-resources:stock}{المخزون} مقسومًا على \omterm{def:g10:freshwater-resources:stock}{الدفق}.
\end{solution}

\begin{solution}{exo:g10:freshwater-resources:2}
نحو \qty{2.5}{\percent} من كل الماء عذب. ومن الماء العذب،
\qty{68.7}{\percent} متجمد في الغطاءات الجليدية والأنهار الجليدية والثلوج الدائمة،
و\qty{30.1}{\percent} تحت الأرض. والأنهار تحوي أقله،
\num{2120}~\unit{km^3}.
\end{solution}

\begin{solution}{exo:g10:freshwater-resources:3}
\omterm{def:g10:freshwater-resources:renewable}{المياه العذبة المتجددة} هي الماء العذب الذي يجدده الهطول كل سنة ويجري عبر منطقة
جريانًا سطحيًا ومياهًا جوفية: إنها حجم في السنة، أي \omterm{def:g10:freshwater-resources:stock}{دفق}. \omterm{def:g10:freshwater-resources:fossil-groundwater}{والمياه الجوفية الأحفورية}
ماء دخل طبقة حاملة للمياه منذ زمن بعيد ولا تجدده التغذية الحالية: فهي مخزن لا يمكن
إلا استنزافه، أي \omterm{def:g10:freshwater-resources:stock}{مخزون}.
\end{solution}

\begin{solution}{exo:g10:freshwater-resources:4}
$91\,000/2120 \approx 43$ مرة.
\end{solution}

\begin{solution}{exo:g10:freshwater-resources:5}
الغلاف الجوي يجدد ماءه أسرع (نحو \num{9.5} أيام)، والمحيط أبطأ
الخزانات المحسوبة (نحو \num{3200} سنة)؛ والمياه الجوفية العميقة أقدم أيضًا.
$3200 \times 365.25/9.5 \approx \num{120000}$
مرة أطول.
\end{solution}

\begin{solution}{exo:g10:freshwater-resources:6}
الداخل: $413\,000 + 73\,000 = \num{486000}$~\unit{km^3} في السنة، إذن
$T = 12\,700/486\,000 \approx \num{0.0261}$ سنة $\approx \num{9.5}$ أيام.
والخارج: $373\,000 + 113\,000 = \num{486000}$~\unit{km^3}، أي القدر نفسه: فالهواء في
حالة ثبات، \omterm{def:g10:freshwater-resources:residence-time}{وزمن المكوث} هو نفسه.
\end{solution}

\begin{solution}{exo:g10:freshwater-resources:7}
$1\,335\,040/413 \approx \num{3230}$ سنة. وفي حالة الثبات يتلقى المحيط
قدر ما يفقده: $373 + 40 = 413$ ألف~\unit{km^3} في السنة،
\omterm{def:g10:freshwater-resources:residence-time}{فزمن المكوث} المحسوب بالمداخل هو نفسه.
\end{solution}

\begin{solution}{exo:g10:freshwater-resources:8}
تتلقى اليابسة \num{113000}~\unit{km^3} في السنة؛ يعود منها \num{73000} إلى
الهواء بالتبخر والنتح ويجري \num{40000} إلى البحر:
$73\,000 + 40\,000 = \num{113000}$. والحصة العائدة إلى الهواء
$73/113 \approx \qty{65}{\percent}$.
\end{solution}

\begin{solution}{exo:g10:freshwater-resources:9}
$273.2 \times 10^6 \times 1233.5 \approx \qty{3.37e11}{m^3}$، أي
\qty{337}{km^3}: قرابة كل الماء الذي كان بحر آرال لا يزال يحويه في
1987.
\end{solution}

\begin{solution}{exo:g10:freshwater-resources:10}
$15.8 \times 0.3048 \approx \qty{4.82}{m}$. وعلى مدى السنوات الخمس والستين من 1950 إلى
2015، $482/65 \approx \qty{7.4}{cm}$ في السنة.
\end{solution}

\begin{solution}{exo:g10:freshwater-resources:11}
$2/0.01 = 200$ سنة. وينخفض \omterm{def:g10:freshwater-resources:stock}{المخزون} الآن بمقدار $0.05 - 0.01 =
\qty{0.04}{km^3}$ في السنة، فيدوم $2/0.04 = 50$ سنة.
\end{solution}

\begin{solution}{exo:g10:freshwater-resources:12}
1960: $1093/71.1 \approx 15$ سنة؛ 2010: $98.1/11.4 \approx 8.6$ سنوات.
يزيل التبخر كل سنة سماكة من الماء تكاد تكون واحدة من كل كيلومتر مربع من
البحيرة؛ ومع تقلص البحيرة، تقلصت مساحتها، وتقلص معها الحجم الكلي المتبخر.
\end{solution}

\begin{solution}{exo:g10:freshwater-resources:13}
(1) الفقدان الصافي $1.2 - 0.2 = \qty{1.0}{m}$ في السنة على
$100 \times 10^6$~\unit{m^2}: أي $10^8$~\unit{m^3}، فيجب أن يرد \qty{0.1}{km^3}
في السنة. (2) ينخفض الوارد إلى \qty{0.05}{km^3}: فتفقد البحيرة
$\qty{0.05}{km^3} = 5 \times 10^7$~\unit{m^3}، أي طبقة سماكتها
$5 \times 10^7/10^8 = \qty{0.5}{m}$. (3) مع انخفاض المنسوب، تنكشف الشواطئ
من الماء وتتقلص مساحة البحيرة: فيقل السطح، ويقل التبخر، حتى يتوازن التبخر
والوارد من جديد.
\end{solution}

\begin{solution}{exo:g10:freshwater-resources:14}
تحوي الطبقة $100 \times 0.20 = \qty{20}{m}$ من الماء، أي
\qty{20000}{mm}: فبمعدل \qty{1}{mm} في السنة، \num{20000} سنة. والماء العميق في
الطبقة الحاملة للمياه في الصحراء الكبرى يتراوح عمره بين \num{200000} سنة ومليون سنة، أي
أكثر بعشر مرات إلى خمسين مرة: فقد غُذّي في أثناء مناخات ماضية أرطب.
\end{solution}

\begin{solution}{exo:g10:freshwater-resources:15}
$337/453\,000 \approx \qty{0.00074}{km}$، أي نحو \qty{0.74}{m} من الماء. وقد انخفض
منسوب الماء \qty{4.8}{m}، أي أكثر بنحو ست مرات ونصف. فالماء لا
يملأ إلا مسام الصخر (ولا ينصرف كله منها): فحين ينخفض منسوب المياه الجوفية مترًا، تكون
طبقة الماء المحرَّر أرق بكثير من متر، وهي هنا نحو $0.74/4.8 \approx \num{0.15}$~m.
\end{solution}

\begin{solution}{pb:g10:freshwater-resources:1}
\textbf{1.} $1093/68\,000 \approx \qty{0.016}{km}$، أي نحو \qty{16}{m}.

\textbf{2.} $1093 \times 10^9 \approx \qty{1.09e12}{m^3}$.

\textbf{3.} \qty{10}{g/L} تساوي \qty{10}{kg/m^3}: $1.093 \times 10^{12} \times 10
\approx \qty{1.1e13}{kg}$، أي نحو \num{11} مليار طن.

\textbf{4.} $1093/71.1 \approx 15$ سنة.

\textbf{5.} لا يخرج منها أي نهر: فلا يغادرها ماؤها إلا بخارًا (أو متسربًا
في الأرض)، بينما يبقى الملح.

\textbf{6.} $71.1/68\,000 \approx \qty{0.00105}{km}$، أي نحو \qty{1.05}{m}.

\textbf{7.} $9.4/68\,000 \approx \qty{0.00014}{km}$، أي نحو \qty{0.14}{m}.

\textbf{8.} $71.1 - 9.4 = \qty{61.7}{km^3}$ في السنة.

\textbf{9.} $79 + 37 = \qty{116}{km^3}$: $61.7/116 \approx \qty{53}{\percent}$،
أي نحو النصف.

\textbf{10.} المداخل $1.0 + 6.2 = \qty{7.2}{km^3}$؛ والمخرج \qty{36.8}{km^3}؛
والتغير الصافي $7.2 - 36.8 = \qty{-29.6}{km^3}$.

\textbf{11.} $(502.7 - 327.2)/5 = \qty{35.1}{km^3}$ في السنة: من رتبة
\qty{29.6}{km^3} في ميزانية 1987 نفسها (والسنوات المفردة تختلف: فقد كان التبخر
\qty{47.9}{km^3} في 1984).

\textbf{12.} $36.8/0.00105 \approx \qty{35000}{km^2}$.

\textbf{13.} $35\,000/68\,000 \approx \qty{52}{\percent}$، أي نحو النصف.

\textbf{14.} تقلصت البحيرة: ونصف السطح يبخر نحو نصف
الماء.

\textbf{15.} $345.6/1093 \approx \qty{32}{\percent}$، أي نحو الثلث.

\textbf{16.} $10 \times 1093/345.6 \approx \qty{31.6}{g/L}$.

\textbf{17.} الملوحة المقيسة أدنى: فقد غادر جزء من الملح
الماء، فتبلور على القاع المكشوف وحملته الريح غبارًا
مالحًا.

\textbf{18.} نحو \qty{1.05}{m} في السنة، وهي السماكة المتبخرة.

\textbf{19.} $7.2/35\,000 \approx \qty{0.0002}{km}$، أي نحو \qty{0.20}{m}.

\textbf{20.} $1.05 - 0.20 \approx 0.85$ (و\num{0.84} بالقيم غير المقرّبة):
\textbf{في 1987 انخفض منسوب بحر آرال بنحو \qty{0.8}{m}، أي قرابة
متر في سنة واحدة}.
\end{solution}
